\documentclass[../main.tex]{subfiles}
\begin{document}

\textit{Chương này trình bày hệ thống đã cài đặt qua ảnh chụp từ ứng dụng đang chạy, sau đó là chiến
lược kiểm thử, kết quả và việc triển khai.}

\section{Môi trường phát triển}
\begin{table}[H]
\centering
\caption{Môi trường phát triển và quy mô dự án}\label{tab:env}
\renewcommand{\arraystretch}{1.2}
\begin{tabular}{|p{5cm}|p{8.7cm}|}
\hline
Hệ điều hành & Ubuntu Linux, Docker Engine với Compose v2 \\ \hline
Frontend & Node.js 22, pnpm 10, Vite, Vue 3, TypeScript \\ \hline
Backend & Python 3.12, FastAPI, SQLAlchemy 2, Alembic \\ \hline
Cơ sở dữ liệu & MySQL 8.4 (production), SQLite (phát triển, kiểm thử) \\ \hline
Mã nguồn & khoảng 13.300 dòng frontend, 3.400 dòng backend (không tính dòng trống) \\ \hline
Mã kiểm thử & khoảng 5.200 dòng frontend, 2.800 dòng backend \\ \hline
Kho mã & 65 commit; 26 component Vue; 51 endpoint API; 10 bảng; 9 migration \\ \hline
\end{tabular}
\end{table}

\section{Các chức năng đã cài đặt}
\subsection{Trang giới thiệu và xác thực}
Khách truy cập trước hết thấy trang giới thiệu với các tính năng, video demo và bảng giá; nút đăng
nhập nằm ở góc trên bên phải và có thể chuyển ngôn ngữ Anh/Việt (Hình~\ref{fig:landing.png}). Video
demo là bản quay màn hình thật: một script điều khiển ứng dụng thật trên trình duyệt và ghi lại, nên
nó chỉ có thể cho thấy những chức năng thực sự chạy. Khi đăng ký chỉ cần tên, email và mật khẩu---không
bao giờ hỏi vai trò (Hình~\ref{fig:login.png}).

\shot{landing.png}{0.95}{Trang giới thiệu}
\shot{login.png}{0.95}{Hộp thoại đăng nhập}

\subsection{Dự án, thành viên và vai trò}
Bảng điều khiển liệt kê các dự án của người dùng kèm vai trò trong từng dự án, cho phép tạo dự án và
quản lý nhóm (Hình~\ref{fig:dashboard.png}). Thêm thành viên qua email sẽ gửi lời mời; người đó chỉ
tham gia sau khi chấp nhận.

\shot{dashboard.png}{0.95}{Bảng điều khiển dự án với vai trò theo từng dự án}

\subsection{Bảng Kanban và dạng bảng}
Bảng có bốn cột. Mỗi thẻ hiển thị mã, tiêu đề, độ ưu tiên, ảnh đại diện người thực hiện, hạn chót
(``còn 9 ngày'', ``trễ 3 ngày''), lý do bị chặn và biểu tượng ngọn lửa cho task trên đường găng. Thứ
tự sắp xếp ưu tiên hạn chót: một task ưu tiên thấp hết hạn hôm qua đứng trên một task khẩn cấp hạn
tháng sau. Cùng dữ liệu đó có thể xem dạng bảng (Hình~\ref{fig:board-kanban.png} và
\ref{fig:board-table.png}); phím \code{b} để chuyển đổi.

\shot{board-kanban.png}{0.95}{Bảng Kanban}
\shot{board-table.png}{0.95}{Dạng bảng}

Mọi thao tác đều có phím tắt (Hình~\ref{fig:shortcuts.png}): $h/j/k/l$ di chuyển, Space đổi trạng
thái, $H/L$ chuyển task giữa các cột, 1--4 đặt độ ưu tiên, $n$ tạo task và $/$ tìm kiếm.

\shot{shortcuts.png}{0.8}{Phím tắt}

\subsection{Chi tiết task: thảo luận, minh chứng và lịch sử}
Mở một task sẽ thấy các trường, phụ thuộc và tiêu chí chấp nhận, luồng thảo luận có @mention và trả
lời một cấp, file đính kèm, và lịch sử thay đổi do server ghi lại (Hình~\ref{fig:task-detail.png}).
Khi chuyển task sang DONE, hệ thống mời người dùng đính kèm minh chứng rằng việc đã xong.

\shot{task-detail.png}{0.95}{Chi tiết task với thảo luận và lịch sử}

\subsection{Scrum}
Bảng Scrum có năm tab. \emph{Sprint} (Hình~\ref{fig:scrum-sprints.png}) liệt kê các sprint đã lên
kế hoạch, đang chạy và đã xong cùng số điểm và mức cam kết; PM tạo, bắt đầu và kết thúc sprint, chọn
nơi chuyển việc chưa xong. \emph{Lập kế hoạch} (Hình~\ref{fig:scrum-planning.png}) chuyển việc giữa
product backlog và sprint. \emph{Burndown} (Hình~\ref{fig:scrum-burndown.png}) vẽ điểm còn lại so với
đường lý tưởng; \emph{Velocity} (Hình~\ref{fig:scrum-velocity.png}) so sánh điểm cam kết và điểm hoàn
thành của từng sprint cùng giá trị trung bình; \emph{Retrospective}
(Hình~\ref{fig:scrum-retro.png}) nhờ AI nêu điều làm tốt, điều cần cải thiện và hành động cụ thể, dựa
trên số liệu của sprint.

\shot{scrum-sprints.png}{0.95}{Bảng Scrum: vòng đời sprint}
\shot{scrum-planning.png}{0.95}{Bảng Scrum: lập kế hoạch backlog}
\shot{scrum-burndown.png}{0.95}{Bảng Scrum: biểu đồ burndown}
\shot{scrum-velocity.png}{0.95}{Bảng Scrum: velocity}
\shot{scrum-retro.png}{0.95}{Bảng Scrum: retrospective do AI viết}

\subsection{Các công cụ AI}
Bốn công cụ AI khác nằm trong menu \emph{Công cụ AI}. Tóm tắt tuần báo cáo tiến độ, rủi ro và ưu
tiên (Hình~\ref{fig:ai-weekly.png}); trình tạo biên bản họp biến ghi chép thô thành chủ đề, việc cần
làm có người phụ trách và độ ưu tiên, cùng các quyết định (Hình~\ref{fig:ai-minutes.png}); giao việc
thông minh hiển thị khối lượng và kỹ năng từng thành viên rồi gợi ý người thực hiện kèm lý do và đánh
giá rủi ro (Hình~\ref{fig:ai-assign.png}); công cụ chia nhỏ mục tiêu tách một mục tiêu thành các task
con có phụ thuộc, có thể chèn vào bảng bằng một thao tác (Hình~\ref{fig:ai-decomposer.png}).

\shot{ai-weekly.png}{0.95}{AI tóm tắt tiến độ tuần}
\shot{ai-minutes.png}{0.95}{AI tạo biên bản họp}
\shot{ai-assign.png}{0.95}{AI phân tích khối lượng và gợi ý giao việc}
\shot{ai-decomposer.png}{0.95}{AI chia nhỏ mục tiêu}

\subsection{Đồ thị phụ thuộc, thông báo và hồ sơ}
Màn hình DAG liệt kê task theo thứ tự topo và đánh dấu đường găng (Hình~\ref{fig:dag.png}). Luồng
thông báo báo các lần được nhắc tên, được trả lời, được giao việc và được mời---không bao giờ báo về
hành động của chính người dùng (Hình~\ref{fig:notifications.png}). Trang hồ sơ cho sửa tên, kỹ năng
(được dùng để gợi ý giao việc) và ảnh đại diện (Hình~\ref{fig:profile.png}).

\shot{dag.png}{0.95}{Thứ tự thực hiện DAG và đường găng}
\shot{notifications.png}{0.95}{Luồng thông báo}
\shot{profile.png}{0.95}{Trang hồ sơ}

\section{Kiểm thử}
\subsection{Chiến lược}
Kiểm thử gồm bốn lớp (Chương~5 phân loại theo mức đơn vị và tích hợp):
\begin{enumerate}
  \item \textbf{Test backend} (pytest, 13 file) gọi API qua test client: mọi quy tắc vai trò và quyền
        (kể cả trường hợp phủ định---thành viên nhận 403, người ngoài nhận 404), vòng trạng thái, quy
        tắc Scrum, lịch sử, migration, kiểm tra an toàn khi khởi động, và chuỗi AI với từng tầng được
        giả lập ở tầng HTTP.
  \item \textbf{Test frontend} (Vitest, 19 file) bao phủ các thuật toán thuần, cả hai store, bộ xử
        lý phím tắt và các component trong DOM giả lập. Một test đọc template của mọi component và
        báo lỗi nếu còn chữ chưa được dịch.
  \item \textbf{Smoke test đầu-cuối}: một script nạp dữ liệu vào instance đang chạy qua API công khai
        và thực hiện 32 bước kiểm tra trên mọi chức năng, kể cả bốn endpoint AI với mô hình thật.
  \item \textbf{Tự động hoá trình duyệt}: Playwright điều khiển ứng dụng thật để quay video demo và
        chụp ảnh trong chương này, đồng thời quét từng màn hình để tìm chữ sai ngôn ngữ.
\end{enumerate}
Bản thân các test cũng được kiểm tra bằng kiểm thử đột biến: cố ý xoá một điều kiện kiểm tra hoặc
đổi một giá trị, và chỉ chấp nhận khi có test báo lỗi. Ví dụ, xoá điều kiện ``sprint phải thuộc cùng
dự án với task'' làm hai test Scrum báo lỗi; xoá bộ lọc chỉ lấy thành viên đã chấp nhận làm test gợi ý
giao việc báo lỗi.

\subsection{Kết quả}
\begin{table}[H]
\centering
\caption{Kết quả kiểm thử}\label{tab:tests}
\renewcommand{\arraystretch}{1.2}
\begin{tabular}{|p{6.5cm}|c|c|}
\hline
\textbf{Bộ kiểm thử} & \textbf{Số test} & \textbf{Kết quả} \\ \hline
Backend (pytest, SQLite) & 185 & đạt toàn bộ \\ \hline
Rào chắn agent (pytest) & 56 & đạt toàn bộ \\ \hline
Frontend (Vitest) & 442 & đạt toàn bộ \\ \hline
Kiểm tra kiểu TypeScript và build production & -- & đạt \\ \hline
Smoke test đầu-cuối trên bản công khai & 32 bước & 0 lỗi \\ \hline
Quét lẫn ngôn ngữ trên 20 màn hình (mỗi ngôn ngữ) & 20 & không phát hiện \\ \hline
\end{tabular}
\end{table}

\subsection{Các lỗi đã phát hiện}
Kiểm thử đã phát hiện nhiều lỗi thật, mỗi lỗi được ghi trong sổ tay phát triển kèm nguyên nhân. Một
số lỗi tiêu biểu:
\begin{longtable}{|p{1.2cm}|p{9cm}|p{3.3cm}|}
\caption{Một số lỗi đã phát hiện và sửa}\label{tab:defects}\\
\hline
\textbf{Mã} & \textbf{Lỗi} & \textbf{Cách phát hiện} \\ \hline
\endfirsthead
\hline
\textbf{Mã} & \textbf{Lỗi} & \textbf{Cách phát hiện} \\ \hline
\endhead
F-32 & Image backend phát hành chứa file \code{.env} của lập trình viên cùng khoá ký phiên và một CSDL. & Kiểm tra image đã build \\ \hline
F-37 & Lệnh triển khai trong tài liệu chép cấu hình và CSDL phát triển đè lên production. & Đọc kỹ lệnh trước khi chạy \\ \hline
F-45 & Mọi file đính kèm và ảnh đại diện hiển thị lỗi: thẻ \code{<img>} không gửi được header xác thực. & Kiểm thử thủ công \\ \hline
F-47 & Khoá ngoại SQLite chưa bao giờ được bật, xoá task để lại dữ liệu mồ côi. & Khi chuyển sang MySQL \\ \hline
F-48 & Có thể xếp task vào sprint của dự án khác. & Khi viết test Scrum \\ \hline
F-49 & Task đã xong vẫn bị ghi ``trễ 10 ngày''. & Xem lại ảnh chụp \\ \hline
F-50 & Công cụ ``AI'' chia nhỏ mục tiêu chưa từng gọi mô hình, chỉ trả về mẫu cố định. & Thử với mô hình thật \\ \hline
F-51 & Gợi ý giao việc có thể chọn người chưa chấp nhận lời mời. & Rà soát mã khi kiểm thử AI \\ \hline
F-53 & Một thay đổi được đề xuất giả lập vai trò PM ngay trên trình duyệt sau khi server từ chối. & Rà soát thay đổi được gửi tới \\ \hline
-- & 21/26 component có chữ chưa dịch, khiến giao diện lẫn hai ngôn ngữ. & Thành viên phản hồi; đã sửa và có test chặn \\ \hline
\end{longtable}

\section{Triển khai}
Cấu hình production chạy bốn container: MySQL, backend FastAPI, nginx cùng frontend đã build, và
\code{cloudflared}. Script triển khai build image, xác minh trong image không có khoá bí mật hay CSDL,
sao lưu CSDL, chạy migration ở một bước riêng, khởi động dịch vụ và kiểm tra instance hoạt động tốt,
ở phiên bản lược đồ mới nhất và đúng tư thế production (không có tài khoản demo, không có CORS
wildcard, khoá bí mật không phải mặc định). Tunnel công khai instance tại
\url{https://koshi.huynhca.com}, nơi nhóm đã kiểm thử với ba tài khoản được tạo sẵn. Tích hợp liên tục
trên GitHub Actions chạy các bộ test frontend, backend và rào chắn mỗi lần push lên nhánh chính và
nhánh phát triển. Vận hành và theo dõi được trình bày ở Mục~5.6.

\end{document}
